\chapter{文本表示：从符号到模型输入}
\label{ch:text}

文本要先转成模型能处理的表示。词袋、TF-IDF与嵌入向量保留的信息不同，
并非一个必然比另一个更高级。先用简单基线理解任务，再决定是否需要序列与上下文模型。
本章先建立数值表示，下一章再讨论不同位置之间怎样交换信息。

\begin{keyidea}[符号编号没有大小距离含义]
Token ID是查表索引，不表示语义的高低远近。词典、分词规则与模型输入必须一致；
需要从当前数据学习的词典和文档频率，应在训练内拟合。
\end{keyidea}

\section{Tokenization、词干与停用词}
\label{sec:text-tokens}
Tokenization把文本切成处理单元（Token），可以是词、子词、字符等。
中文不必按空格切分，词干提取也不是所有语言和模型都需要的步骤。

\begin{figure}[htbp]
\centering\begin{tikzpicture}[node distance=8mm]
\node[box,text width=25mm] (text) {文本\\我爱猫};
\node[box,text width=30mm,right=of text] (tok) {Token序列\\我 / 爱 / 猫};
\node[box,text width=23mm,right=of tok] (ids) {ID序列\\1 / 2 / 3};
\node[box,text width=30mm,right=of ids] (emb) {嵌入查表\\$3\times d$};
\draw[flow] (text)--(tok);\draw[flow] (tok)--(ids);\draw[flow] (ids)--(emb);
\end{tikzpicture}

\caption{一个人工分词与词典例子，词典含PAD=0、我=1、爱=2、猫=3。实际系统可采用不同切分。}
\par\smallskip
\readfigure{三项ID查三行嵌入，得到三条向量。ID为3不表示“猫”是“我”的三倍。
验证与上线必须沿用同一词典及未知Token处理规则。}
\end{figure}

\begin{center}
\begin{tabularx}{\textwidth}{P{35mm}Y}
\toprule
处理 & 作用与限制\\
\midrule
词干提取 Stemming & 按规则缩减词形，输出未必是语言学词根或完整单词。\\
词形还原 Lemmatization & 试图还原词元，可能依赖词性或语言资源。\\
停用词 Stop words & 某些任务过滤高频功能词；否定、语气等也可能是关键信息，不能默认全删。\\
\bottomrule
\end{tabularx}
\end{center}

\quick{“不喜欢”中的“不”可能决定情感标签；文本清洗必须与任务语义一致。}

\section{词袋与TF-IDF：可解释的稀疏基线}
\label{sec:text-tfidf}
词袋（Bag of Words）把文档表示为词典上的计数或出现指示向量。
它忽略词序：“猫追狗”与“狗追猫”在相同分词下可以具有相同词袋，但含义不同。
指示值0/1不是一个“二元分类器”，只是特征数值。

\begin{center}
\begin{tabular}{lrrrr}
\toprule
文档/词典 & 猫 & 狗 & 看 & 跑\\
\midrule
D1：猫 看 猫 & 2 & 0 & 1 & 0\\
D2：狗 看 & 0 & 1 & 1 & 0\\
D3：猫 跑 & 1 & 0 & 0 & 1\\
文档频率$df$ & 2 & 1 & 2 & 1\\
\bottomrule
\end{tabular}
\end{center}

TF-IDF组合词频与逆文档频率。本例使用原始计数$tf$、自然对数和平滑式
\[
 idf(t)=\log\frac{1+N}{1+df(t)}+1,\qquad v_{d,t}=tf(t,d)\,idf(t),
\]
其中$N=3$，$df$按出现该词的\textbf{文档数}计，不是总次数。
“猫”和“看”的$idf=1+\log(4/3)\approx1.288$，另外两项为$1+\log2\approx1.693$。
若再作L2归一化，D1变为$[2/\sqrt5,0,1/\sqrt5,0]$。
输入上述计数矩阵时，平滑、原始计数与L2归一化口径对应常见\texttt{TfidfTransformer}设置；
这里不假定默认分词器会按单个汉字切分。
若一篇新文档没有任何词落在固定词典内，它会形成全零向量，不能直接除以0来归一化。

不同实现有不同TF、平滑与归一化约定，应报告完整设置。
常见词权重降低不代表它没有语义价值，稀有词权重高也不证明它具有预测或因果作用。

\quick{词典与IDF从训练数据学习，随后固定用于验证、测试和新文档。}

\section{嵌入与Word2Vec：学习稠密向量}
\label{sec:text-embedding}
嵌入（Embedding）用矩阵$E\in\R^{V\times d}$为$V$个Token存储$d$维向量。
长度$T$的ID序列查表得到$T\times d$矩阵，批量后常为$B\times T\times d$。
仅嵌入矩阵就有$Vd$个参数；PAD、未知词和是否冻结按实现设计。

Word2Vec通过局部上下文预测等目标学习词向量，典型形式包括CBOW与Skip-gram。
上下文相似可能产生向量相近，但相似性和类比关系是数据与训练的结果，
不是对任意词都成立的语义定理。词向量还可能保留语料中的偏差。

\begin{center}
\begin{tabularx}{\textwidth}{P{35mm}YY}
\toprule
表示 & 同一Token在不同句子中 & 常见考虑\\
\midrule
词袋/TF-IDF & 对应固定词典列。 & 简单、稀疏；通常忽略顺序。\\
静态词嵌入 & 查同一行向量。 & 稠密；多义词共享初始表示。\\
上下文表示 & 经序列/注意力层后可随上下文变化。 & 依结构处理次序与可见范围；训练和使用成本通常更高。\\
\bottomrule
\end{tabularx}
\end{center}

预训练嵌入可冻结后作为特征，也可以随任务微调。预处理应与预训练模型约定一致，
不能任意删词、改变Token划分后仍假定原模型语义不变。

\begingroup\small
\noindent\textbf{参考：}Manning等，\textit{Introduction to Information Retrieval}（2008）第6章；
Mikolov等（2013）的Word2Vec工作；scikit-learn Text feature extraction文档。
具体TF-IDF口径见\href{https://scikit-learn.org/stable/modules/feature_extraction.html#tfidf-term-weighting}{词频加权与归一化说明}。
本章表格与流程图独立设计。
\par\endgroup
